\documentclass[12pt, a4paper]{article}

% ── Codificación e idioma ──────────────────────────────────────────────────────
\usepackage[utf8]{inputenc}
\usepackage[T1]{fontenc}
\usepackage[spanish, es-tabla]{babel}

% ── Márgenes y geometría ───────────────────────────────────────────────────────
\usepackage[top=2.5cm, bottom=2.5cm, left=3cm, right=3cm]{geometry}

% ── Tipografía y espaciado ─────────────────────────────────────────────────────
\usepackage{lmodern}
\usepackage{setspace}
\onehalfspacing

% ── Matemáticas y símbolos ─────────────────────────────────────────────────────
\usepackage{amsmath}
\usepackage{amssymb}

% ── Listas mejoradas ───────────────────────────────────────────────────────────
\usepackage{enumitem}

% ── Tablas y gráficos ──────────────────────────────────────────────────────────
\usepackage{booktabs}
\usepackage{array}
\usepackage{longtable}
\usepackage{graphicx}

% ── Hipervínculos ──────────────────────────────────────────────────────────────
\usepackage{hyperref}
\hypersetup{
    colorlinks=true,
    linkcolor=blue,
    urlcolor=blue,
    citecolor=blue,
    pdftitle={De Data Lakes a Data Mesh: Arquitecturas Distribuidas},
    pdfauthor={Rubén Prieto García, Oscar},
}

% ── Cabeceras y pies de página ─────────────────────────────────────────────────
\usepackage{fancyhdr}
\pagestyle{fancy}
\fancyhf{}
\rhead{\small \textit{Arquitectura Data Mesh}}
\lhead{\small \textit{Sección \thesection}}
\cfoot{\thepage}

% ── Párrafos ───────────────────────────────────────────────────────────────────
\setlength{\parskip}{6pt}
\setlength{\parindent}{0pt}

% ───────────────────────────────────────────────────────────────────────────────
\begin{document}

% ── Portada ────────────────────────────────────────────────────────────────────
\begin{titlepage}
    \centering
    \vspace*{2cm}
    
    {\Large \textbf{Universidad de Jaén (UJA)}\par}
    \vspace{0.5cm}
    {\large Máster en Ingeniería Informática\par}
    
    \vspace{3cm}
    
    {\Huge \textbf{De Data Lakes Monolíticos a Data Mesh}\par}
    \vspace{0.5cm}
    {\Large \textit{Arquitecturas Distribuidas y Gobernanza Federada en Entornos Cloud}\par}
    
    \vspace{3cm}
    
    {\large \textbf{Autores:}\par}
    {\normalsize Rubén Prieto García\par}
    {\normalsize Oscar [Apellidos de Oscar]\par}
    
    \vspace{2cm}
    
    {\large \textbf{Línea temática:}\par}
    {\normalsize Big Data en la nube: Arquitecturas y modelos de implementación\par}
    
    \vfill
    
    {\normalsize \today\par}
\end{titlepage}

\tableofcontents
\newpage

% =============================================================================
\section{Introducción y Contexto del Problema}
% =============================================================================

La evolución de los sistemas de información corporativos ha consolidado a los datos como el activo subyacente más crítico para la toma de decisiones estratégicas. Sin embargo, el aumento exponencial en el volumen, velocidad y variedad de la información ha trasladado el desafío de la ingeniería de datos: ya no se trata de un problema de capacidad de almacenamiento o cómputo bruto, sino de un problema de escalabilidad organizativa, gobernanza y arquitectura distribuida. Este proyecto propone una revisión crítica de los modelos centralizados y diseña una solución técnica basada en topologías federadas.

% -----------------------------------------------------------------------------
\subsection{El Cuello de Botella Organizativo}
% -----------------------------------------------------------------------------

Durante la última década, la industria tecnológica ha abordado la gestión de datos a gran escala mediante el patrón arquitectónico del \textit{Data Lake} (lago de datos). Esta aproximación centralizada supuso un avance significativo frente a los tradicionales \textit{Data Warehouses} (sistemas de almacenamiento estructurado optimizados para consultas históricas predefinidas), al permitir la ingesta masiva de datos en repositorios unificados de bajo coste, separando funcionalmente la capa de almacenamiento de la de computación.

No obstante, a medida que las organizaciones escalan en complejidad, este paradigma centralizado ha revelado limitaciones sistémicas severas. El crecimiento de los datos ha generado una disfunción conocida como el \textbf{``cuello de botella del equipo central''}. En las arquitecturas monolíticas actuales, un único equipo de ingenieros de datos asume la responsabilidad de extraer, limpiar, transformar y servir la información de todos los departamentos de la empresa. Esta centralización provoca fallos críticos:

\begin{itemize}
    \item \textbf{Pérdida de Contexto Semántico:} El equipo central carece del conocimiento de dominio necesario para entender la lógica de negocio, convirtiendo el \textit{Data Lake} en un pantano de datos (\textit{Data Swamp} - repositorios caóticos e inmanejables) con información no catalogada ni aprovechable por la organización (\textit{Dark Data}).
    \item \textbf{Latencia Operativa (\textit{Time-to-Market}):} Los consumidores (analistas o sistemas analíticos) sufren tiempos de espera inasumibles para acceder a datasets limpios debido a la sobrecarga del nodo central de procesamiento humano.
\end{itemize}

% -----------------------------------------------------------------------------
\subsection{De Data Lakes a Arquitecturas Distribuidas}
% -----------------------------------------------------------------------------

En las arquitecturas analíticas convencionales, la responsabilidad técnica y el conocimiento del dominio se encuentran inherentemente desacoplados. Los productores vuelcan datos en bruto delegando el procesamiento posterior. Desde la perspectiva de la computación distribuida, esta topología genera un cuello de botella de recursos compartidos. 

Cualquier cambio en el esquema de origen de un departamento requiere la intervención manual del equipo central, bloqueando los despliegues de otros dominios y rompiendo los principios de integración continua. El problema fundamental no es la capacidad tecnológica de procesar Big Data, sino la incapacidad de escalar la arquitectura técnica sin distribuir la computación y la propiedad del dato hacia la periferia (los dominios).

% -----------------------------------------------------------------------------
\subsection{Objetivos del Proyecto}
% -----------------------------------------------------------------------------

Para resolver las limitaciones sistémicas descritas, este proyecto establece los siguientes objetivos técnicos y arquitectónicos:

\begin{itemize}
    \item \textbf{Objetivo Principal:} Diseñar e implementar una arquitectura de datos distribuida basada en el paradigma \textit{Data Mesh}, demostrando la viabilidad técnica de descentralizar la propiedad analítica mediante la orquestación de clústeres independientes e interoperables.
    \item \textbf{Simulación de Interoperabilidad Federada:} Modelar la interacción asíncrona entre dos nodos computacionales autónomos (Dominio Productor y Consumidor) mediante la creación de ``Productos de Datos'' estandarizados.
    \item \textbf{Implementación de Gobernanza Computacional:} Establecer un plano de control unificado que aplique automáticamente políticas de seguridad y control de acceso basado en roles (RBAC, que restringe el acceso a la red según la función del usuario en la empresa) sin comprometer la autonomía local.
\end{itemize}

% -----------------------------------------------------------------------------
\subsection{Alcance y Estrategia FinOps (Local-First)}
% -----------------------------------------------------------------------------

Para garantizar la rentabilidad, escalabilidad y reproducibilidad del proyecto, la metodología de desarrollo se acota bajo las siguientes premisas:

\begin{enumerate}
    \item \textbf{Estrategia ``Local-First, Cloud-Ready'' (FinOps):} Para maximizar la rentabilidad durante el desarrollo y demostrar prácticas de optimización de costes, la Prueba de Concepto (PoC) no se desplegará sobre infraestructura pública de pago. Se implementará un entorno de simulación local mediante contenerización (Docker), utilizando herramientas 100\% compatibles con las APIs de Amazon Web Services (ej. LocalStack, MinIO, Apache Spark).
    \item \textbf{Infraestructura como Código (IaC) Agnóstica:} La provisión del plano de autoservicio se automatizará mediante lenguajes declarativos (Terraform). El código estará diseñado para que la transición hacia los servicios reales de AWS requiera únicamente la reconfiguración de los \textit{endpoints}, garantizando la portabilidad del sistema.
    \item \textbf{Aproximación Híbrida (Patrón Estrangulador):} Este patrón arquitectónico consiste en migrar gradualmente un sistema heredado reemplazando progresivamente sus funcionalidades. Bajo esta premisa, no se migrará un sistema completo. Se interceptará y refactorizará un flujo de datos estático proveniente de un \textit{Data Lake} heredado simulado, transformándolo en un producto de datos aislado, dejando el resto del monolito intacto.
    \item \textbf{Limitación Analítica:} El proyecto se centra estrictamente en la arquitectura de infraestructura. No se desarrollarán algoritmos de Inteligencia Artificial para la explotación de los datos en esta fase.
\end{enumerate}

% =============================================================================
\section{Fundamentos y Estado del Arte: De Topologías Centralizadas a Data Mesh}
% =============================================================================

El cambio de paradigma hacia topologías federadas trasciende la reestructuración organizativa; constituye la aplicación estricta de los principios de diseño de sistemas distribuidos a la ingeniería de datos. Para comprender la inviabilidad operativa de las arquitecturas centralizadas, es imperativo abstraer el modelo analítico empresarial y evaluarlo bajo las leyes del rendimiento computacional, la tolerancia a particiones y el coste total de propiedad (TCO).

% -----------------------------------------------------------------------------
\subsection{Modelado Matemático de la Centralización}
% -----------------------------------------------------------------------------

En un \textit{Data Lake} monolítico, la infraestructura central y su equipo de ingeniería actúan como un único nodo de procesamiento que recibe peticiones concurrentes (ingesta, transformación, resolución de incidencias) de múltiples dominios. Este comportamiento se modela mediante la \textbf{Teoría de Colas}, específicamente bajo un modelo $M/M/1$ (llegadas markovianas, tiempos de servicio markovianos, un único servidor lógico).

El tiempo de respuesta promedio del sistema ($T$) está determinado por la tasa de llegada de peticiones ($\lambda$) y la tasa de procesamiento del equipo/nodo central ($\mu$):

\begin{equation}
    T = \frac{1}{\mu - \lambda}
\end{equation}

A medida que la empresa se digitaliza, $\lambda$ crece exponencialmente. Cuando la carga de trabajo ($\lambda$) se aproxima a la capacidad máxima de procesamiento ($\mu$), el tiempo de espera tiende a infinito. Este cuello de botella matemático es insalvable sin paralelización.

El impacto global de esta limitación se cuantifica mediante la \textbf{Ley de Amdahl}, que establece que la aceleración máxima teórica ($S$) de un sistema está topada por la fracción de las operaciones que deben ejecutarse secuencialmente ($1 - p$):

\begin{equation}
    S = \frac{1}{(1 - p) + \frac{p}{N}}
\end{equation}

En los modelos centralizados, la limpieza, catalogación y control de calidad representan el factor secuencial ($1 - p$). Independientemente de los recursos de ingesta o consumo ($N$) que se añadan, el \textit{Time-to-Market} analítico de la organización siempre estará limitado asintóticamente por el nodo central. La única solución eficiente es maximizar $p$ distribuyendo la carga computacional hacia los orígenes de los datos.

% -----------------------------------------------------------------------------
\subsection{Topologías Analíticas y el Teorema CAP}
% -----------------------------------------------------------------------------

La distribución de la carga exige abandonar la topología en estrella (\textit{Hub-and-Spoke}) en favor de una \textbf{Topología de Malla Parcialmente Conectada}. En este entorno, los dominios operan como nodos autónomos que se comunican de forma asíncrona mediante el paso de mensajes (eventos o \textit{datasets} inmutables).

Al transicionar hacia un sistema verdaderamente distribuido, la arquitectura queda sujeta al \textbf{Teorema CAP} y su extensión \textbf{PACELC}. En un \textit{Data Mesh}, cuando ocurre una partición de red ($P$), el sistema prioriza la Disponibilidad ($A$) sobre la Consistencia estricta ($C$). Los nodos consumidores pueden seguir operando sobre los productos de datos cacheados o replicados localmente, asumiendo una \textit{Consistencia Eventual}. Durante el funcionamiento normal ($E$), la arquitectura sacrifica la consistencia inmediata en favor de una menor Latencia ($L$), un compromiso óptimo para el procesamiento analítico masivo que no requiere transaccionalidad ACID en tiempo real.

% -----------------------------------------------------------------------------
\subsection{El Paradigma Data Mesh (Los 4 Pilares)}
% -----------------------------------------------------------------------------

Para gobernar esta topología de red distribuida, el patrón \textit{Data Mesh} define cuatro pilares sociotécnicos que actúan como protocolos de interoperabilidad:

\begin{enumerate}
    \item \textbf{Propiedad Orientada a Dominios (Aislamiento de fallos):} Los equipos generadores asumen la responsabilidad integral de la limpieza y exposición de sus datos. Si un nodo productivo falla, el impacto se aísla localmente, previniendo fallos en cascada en toda la red empresarial.
    \item \textbf{Datos como Producto (Interfaces estandarizadas):} El dato no se expone como un archivo crudo, sino encapsulado bajo un ``Contrato de Datos'' (SLA). Este contrato actúa como una API rigurosa que garantiza calidad, descubribilidad y esquemas tipados.
    \item \textbf{Plataforma de Datos de Autoservicio (Plano de control):} El equipo de infraestructura (\textit{Platform Engineering}) abstrae la complejidad proveyendo plantillas agnósticas de Infraestructura como Código (IaC). Esto permite a los dominios instanciar clústeres efímeros de forma autónoma.
    \item \textbf{Gobernanza Computacional Federada (Consenso global):} Se implementa un registro centralizado de políticas y metadatos. Las reglas de seguridad (RBAC/ABAC) y enmascaramiento se definen globalmente pero se ejecutan localmente en los motores de cómputo de cada dominio.
\end{enumerate}

% -----------------------------------------------------------------------------
\subsection{Modelo de Rentabilidad y Coste Total de Propiedad (TCO)}
% -----------------------------------------------------------------------------

La adopción de una arquitectura distribuida conlleva un cambio radical en la estructura del \textit{Total Cost of Ownership} (TCO). Desde el punto de vista financiero y de retorno de inversión (ROI), la rentabilidad del \textit{Data Mesh} se justifica al invertir el balance de costes de infraestructura \textit{Cloud}.

En un modelo monolítico, el mayor sobrecoste proviene de la \textbf{computación ociosa}. Mantener clústeres de procesamiento masivos (\textit{always-on}) para absorber picos de carga centralizada genera un gasto ineficiente. La plataforma de autoservicio del \textit{Data Mesh} mitiga este problema instanciando micro-clústeres analíticos efímeros bajo un modelo \textit{Serverless} (computación sin servidor donde el proveedor asigna recursos dinámicamente y cobra solo por el tiempo exacto de uso), lo que incurre en costes únicamente durante los segundos de ejecución del dominio.

El factor de contrapeso es el \textbf{coste de transferencia de red (\textit{Network Egress})}. Al distribuir los datos, la interoperabilidad entre dominios genera tráfico interno facturable. Sin embargo, el análisis empírico demuestra que el ahorro derivado de la reducción de cuellos de botella organizativos y la drástica disminución del \textit{Time-to-Market} para iniciativas de negocio supera con creces el incremento de los costes de red, garantizando un ROI positivo en entornos corporativos complejos.

% =============================================================================
\section{Diseño de Arquitectura de Referencia}
% =============================================================================

El diseño de la arquitectura de referencia tiene como propósito materializar los principios abstractos del \textit{Data Mesh} en una topología de red gobernable, escalable y tecnológicamente agnóstica. Para garantizar la viabilidad financiera del proyecto sin comprometer su despliegue futuro en producción, el diseño establece una abstracción estricta entre la lógica de negocio distribuida y la infraestructura subyacente.

% -----------------------------------------------------------------------------
\subsection{Equivalencia Tecnológica Híbrida: Estrategia Local-First}
% -----------------------------------------------------------------------------

La orquestación de la Prueba de Concepto (PoC) se ejecuta sobre un entorno local contenerizado orientado a la nube. Para garantizar la inmutabilidad de la Infraestructura como Código (IaC) y permitir una transición directa a Amazon Web Services (AWS) mediante el simple cambio de variables de entorno, se ha diseñado la siguiente matriz de equivalencia tecnológica (ver Tabla \ref{tab:equivalencia}).

\newpage

\begin{table}[h]
    \centering
    \renewcommand{\arraystretch}{1.3}
    \begin{tabular}{>{\raggedright\arraybackslash}p{3cm} >{\raggedright\arraybackslash}p{3.5cm} >{\raggedright\arraybackslash}p{3.5cm} >{\raggedright\arraybackslash}p{4cm}}
        \toprule
        \textbf{Componente Data Mesh} & \textbf{Entorno de Emulación (PoC)} & \textbf{Servicio Gestionado (AWS)} & \textbf{Justificación Técnica} \\
        \midrule
        Almacenamiento Distribuido & MinIO (Docker) & Amazon S3 & Compatibilidad 100\% con API S3. Permite aislamiento de \textit{buckets} por dominio. \\
        Computación y Transformación & Apache Spark (Contenedor efímero) & AWS Glue / Amazon EMR & Ejecución de procesamiento distribuido desacoplado del almacenamiento. \\
        Catálogo Federado & Hive Metastore / LocalStack & AWS Glue Data Catalog & Mantenimiento del \textit{Control Plane} para el registro de productos de datos. \\
        Consumo Analítico & Trino (Docker) & Amazon Athena & Motor de consultas SQL distribuidas sobre esquemas federados. \\
        \bottomrule
    \end{tabular}
    \caption{Matriz de equivalencia tecnológica entre el entorno local y la arquitectura Cloud nativa.}
    \label{tab:equivalencia}
\end{table}

Esta arquitectura híbrida demuestra que los principios de computación distribuida son independientes del proveedor final, maximizando el Retorno de Inversión (ROI) durante la fase de desarrollo.

% -----------------------------------------------------------------------------
\subsection{Topología de Red y Separación de Planos}
% -----------------------------------------------------------------------------

La arquitectura de red rompe el monolito central separando funcionalmente el sistema en dos planos superpuestos pero tecnológicamente aislados:

\begin{itemize}
    \item \textbf{Plano de Datos (\textit{Data Plane}):} Es la topología \textit{Peer-to-Peer} analítica real. Los datos crudos fluyen hacia el \textit{bucket} del Dominio Productor. Tras ser procesados por los clústeres de cómputo efímeros locales, se exponen como productos de datos. El Dominio Consumidor lee estos datos directamente del almacenamiento del productor, sin que la carga útil (el \textit{payload}) pase nunca por un servidor central.
    \item \textbf{Plano de Control (\textit{Control Plane}):} Es la plataforma centralizada de autoservicio. No almacena datos analíticos, sino metadatos, esquemas JSON y políticas de acceso. Actúa como el enrutador lógico que permite el descubrimiento de productos en la red.
\end{itemize}

\begin{figure}[h]
    \centering
    % \includegraphics[width=0.9\textwidth]{img/diagrama_arquitectura.png}
    \vspace{6cm} % Espacio reservado para insertar el diagrama de red real
    \caption{Diagrama C4 (Nivel 2 - Contenedores) ilustrando la separación entre el Plano de Control y los dominios distribuidos en el Plano de Datos.}
    \label{fig:arquitectura_red}
\end{figure}

% -----------------------------------------------------------------------------
\subsection{Políticas de Seguridad y Gobernanza (DevSecOps)}
% -----------------------------------------------------------------------------

En un entorno descentralizado, la seguridad no puede depender del perímetro de red, requiriendo un enfoque \textit{Zero Trust} (Confianza Cero, modelo de seguridad que asume que ninguna entidad es confiable por defecto y exige verificación continua) aplicado directamente a los datos. El diseño de seguridad de la arquitectura impone las siguientes restricciones fundamentales:

\begin{enumerate}
    \item \textbf{Control de Acceso Basado en Atributos y Roles (ABAC/RBAC):} La autorización para consumir un producto de datos no se gestiona a nivel de infraestructura, sino a nivel de metadato en el Plano de Control (simulando las capacidades de \textit{AWS Lake Formation}). Un dominio consumidor debe solicitar explícitamente el acceso (\textit{grant}) a un producto específico.
    \item \textbf{Aislamiento Criptográfico:} Los datos en reposo mantienen un aislamiento lógico por dominios. En una transición a AWS, esto se traduce en la provisión de claves asimétricas dedicadas mediante AWS KMS (\textit{Key Management Service}) para cada producto de datos.
    \item \textbf{Enmascaramiento de PII (\textit{Personally Identifiable Information}):} Los contratos de datos definen automáticamente qué columnas contienen información sensible (ej. correos electrónicos, datos financieros). Las políticas del catálogo federado aplican técnicas de ofuscación de datos dinámicos (\textit{Dynamic Data Masking}) en tiempo de lectura, garantizando que los motores de consulta (Trino/Athena) devuelvan hashes irreversibles a los consumidores sin los privilegios necesarios, cumpliendo con los estándares de \textit{Compliance} y privacidad corporativa.
\end{enumerate}

% -----------------------------------------------------------------------------
\subsection{Gestión de Ciclo de Vida: Versionado y Evolución de Esquemas}
% -----------------------------------------------------------------------------

En una arquitectura \textit{Peer-to-Peer}, la alta dependencia entre dominios requiere una gestión estricta de la evolución de los esquemas de datos (\textit{Schema Evolution}). Si un dominio productor (ej. Ventas) modifica de forma no retrocompatible la estructura de sus datos (como eliminar una columna crítica o alterar un tipo de dato), podría corromper de manera inmediata los flujos de trabajo del dominio consumidor (ej. Marketing).

Para mitigar este riesgo, la arquitectura implementa los Contratos de Datos como interfaces inmutables. Cuando es necesario realizar un cambio disruptivo (\textit{breaking change}), el dominio productor debe desplegar una nueva versión mayor del contrato (ej. \texttt{ventas\_v2}) mientras mantiene activa la versión anterior (\texttt{ventas\_v1}) durante un periodo de deprecación acordado. Esto permite a los consumidores migrar sus canales de ingesta de forma controlada, garantizando la estabilidad operativa del \textit{Data Mesh}.

% =============================================================================
\section{Implementación de la PoC, Auditoría y DevSecOps}
% =============================================================================

La prueba de concepto materializa la transición desde una topología centralizada hacia un ecosistema \textit{Data Mesh}. Siguiendo las directrices de Infraestructura como Código (IaC) y el Patrón Estrangulador, la implementación se divide en la orquestación del plano de control y la posterior interacción asíncrona entre dominios aislados. El código fuente íntegro (scripts de aprovisionamiento, transformaciones y esquemas JSON) se encuentra documentado en los Anexos de esta memoria.

% -----------------------------------------------------------------------------
\subsection{Despliegue de la Infraestructura Heredada (El Legado)}
% -----------------------------------------------------------------------------

El punto de partida técnico requiere la simulación del cuello de botella analítico. Mediante \textit{Terraform} y contenedores \textit{Docker} (ver Anexo I), se aprovisiona un clúster local de almacenamiento de objetos (MinIO) que actúa como el \textit{Data Lake} monolítico original (\texttt{monolithic-datalake-legacy}). 

En este repositorio central se ingiere un \textit{dataset} crudo y desestructurado que simula transacciones históricas mezcladas con perfiles de clientes. Este entorno carece intencionadamente de catálogos de metadatos y de políticas de control de acceso granulares, reflejando el problema de \textit{Dark Data} descrito en los fundamentos teóricos.

% -----------------------------------------------------------------------------
\subsection{Implementación de la Plataforma de Autoservicio}
% -----------------------------------------------------------------------------

Para dotar a los dominios de autonomía, el equipo de plataforma despliega el Plano de Control de forma declarativa. Este plano está compuesto por dos elementos críticos:

\begin{enumerate}
    \item \textbf{El Catálogo Federado:} Se instancia un \textit{Hive Metastore} (emulando \textit{AWS Glue Data Catalog}) encargado de registrar unívocamente la ubicación y el esquema de los nuevos productos de datos que se generen en la red.
    \item \textbf{El Motor de Políticas:} Se configura un gestor centralizado de permisos (emulando \textit{AWS Lake Formation}) que interceptará cualquier petición de lectura entre dominios para evaluar los privilegios de acceso (RBAC).
\end{enumerate}

La provisión de este entorno se realiza en segundos, demostrando la agilidad del pilar de autoservicio para levantar infraestructuras bajo demanda.

% -----------------------------------------------------------------------------
\subsection{Intervención del Dominio Productor (Ventas)}
% -----------------------------------------------------------------------------

El Dominio de Ventas inicia la aplicación del Patrón Estrangulador asumiendo la propiedad sobre su subconjunto de datos. La intervención se ejecuta en tres fases automatizadas:

\begin{enumerate}
    \item \textbf{Extracción Aislada:} Se instancia un clúster efímero de \textit{Apache Spark} (aislado computacionalmente del resto de la organización). Este motor realiza la ingesta directa desde el \textit{Data Lake} legado, evitando depender del equipo central de TI.
    \item \textbf{Limpieza y Transformación (Pipeline ETL):} El \textit{script} de Spark aplica reglas de validación de calidad de datos (eliminación de duplicados, normalización de divisas) y tipado estricto.
    \item \textbf{Publicación del Contrato:} El dato transformado se persiste en un \textit{bucket} físico propiedad exclusiva de Ventas (\texttt{sales-domain-data-product}). Inmediatamente, se envía una llamada a la API del Catálogo Federado registrando el esquema de datos (Contrato) y marcando columnas sensibles (ej. \textit{email} o \textit{tarjeta de crédito}) bajo una etiqueta de restricción (\textit{PII-Tag}).
\end{enumerate}

% -----------------------------------------------------------------------------
\subsection{Consumo y Validación (Marketing)}
% -----------------------------------------------------------------------------

El ciclo analítico se cierra cuando el Dominio Consumidor (Marketing) requiere cruzar los datos de ventas verificadas con sus propias métricas. 

En lugar de construir canales de copia de datos, Marketing despliega su propio motor analítico interactivo (\textit{Trino}). Mediante consultas SQL federadas, Trino interroga al Catálogo Central para descubrir dónde reside el producto ``Ventas Verificadas'' y confirmar su esquema. Una vez resuelta la dirección, el motor ejecuta la lectura de forma \textit{Peer-to-Peer} (P2P) contra el \textit{bucket} de Ventas, garantizando que el cómputo penaliza únicamente los recursos de Marketing, protegiendo así el rendimiento de la base de datos de Ventas.

% -----------------------------------------------------------------------------
\subsection{Auditoría de Seguridad y Trazabilidad (DevSecOps)}
% -----------------------------------------------------------------------------

La eficacia de la gobernanza federada se audita sometiendo el sistema a un test de intrusión lógica. Se ejecutan dos escenarios de consulta desde el Dominio de Marketing:

\begin{itemize}
    \item \textbf{Escenario 1 (Acceso Autorizado):} Una consulta SQL solicita agregaciones de ventas por zona geográfica. El motor de políticas evalúa el rol de Marketing, permite la conexión P2P y devuelve los resultados con una latencia inferior a 500 ms.
    \item \textbf{Escenario 2 (Bloqueo y Enmascaramiento):} Una consulta SQL intenta extraer la columna cruda de ``Correos Electrónicos''. El Plano de Control detecta la \textit{PII-Tag} en el Contrato de Datos. Al verificar que Marketing no posee el nivel de autorización requerido, el sistema aplica enmascaramiento dinámico, devolviendo \textit{hashes} ofuscados en lugar del dato real.
\end{itemize}

Ambas interacciones quedan registradas de forma inmutable en los \textit{logs} de auditoría del sistema (ver evidencias en Anexo II), garantizando el cumplimiento normativo (Compliance) de la infraestructura sin sacrificar la agilidad operativa de los departamentos.

% =============================================================================
\section{Conclusiones y Hoja de Ruta Tecnológica}
% =============================================================================

La ejecución de este proyecto ha demostrado empíricamente que la limitación fundamental en el procesamiento de Big Data a escala corporativa no reside en la capacidad de cómputo, sino en la topología organizativa y de red. La transición desde un modelo centralizado hacia una arquitectura distribuida no es solo una mejora de rendimiento, sino una necesidad de supervivencia operativa.

% -----------------------------------------------------------------------------
\subsection{Conclusiones Arquitectónicas y de Rentabilidad}
% -----------------------------------------------------------------------------

La evaluación de la Prueba de Concepto (PoC) arroja las siguientes conclusiones verificables:

\begin{enumerate}
    \item \textbf{Viabilidad del Patrón Estrangulador:} La interceptación del monolito legado mediante clústeres aislados de \textit{Apache Spark} ha validado que es posible modernizar infraestructuras críticas sin incurrir en tiempos de inactividad (\textit{Downtime}) ni refactorizaciones masivas. El dominio de Ventas logró autonomía total sobre su ciclo de vida de datos.
    \item \textbf{Eficacia del Aislamiento Zero Trust:} La implementación del Plano de Control y la separación física del Plano de Datos garantizaron el cumplimiento de normativas de privacidad. La auditoría demostró que las consultas P2P del dominio de Marketing fueron gobernadas y enmascaradas en milisegundos sin sobrecargar un servidor central.
    \item \textbf{Optimización FinOps y ROI Técnico:} La estrategia \textit{Local-First, Cloud-Ready} demostró un rigor arquitectónico superior. Al contenerizar la infraestructura manteniendo una equivalencia estricta con las APIs de AWS (S3, Glue, Athena, Lake Formation), el equipo ha construido un activo desplegable en producción con coste de desarrollo cero. Esta portabilidad inmutable mediante \textit{Terraform} asegura un Retorno de Inversión (ROI) altamente competitivo para cualquier comité de negocio o fondo de inversión tecnológica.
\end{enumerate}

% -----------------------------------------------------------------------------
\subsection{Hoja de Ruta: Observabilidad Distribuida y Linaje de Datos}
% -----------------------------------------------------------------------------

Dado el éxito de la prueba de concepto, el siguiente paso crítico en la hoja de ruta de adopción empresarial es la implementación de un plano de observabilidad transversal. La descentralización arquitectónica inherente al \textit{Data Mesh} introduce el desafío operativo del diagnóstico de fallos: cuando un dato de mala calidad llega a un consumidor, identificar si la anomalía se originó en la ingesta inicial, en una transformación intermedia o en el transporte de red requiere una trazabilidad absoluta.

Para resolver este desafío, las futuras iteraciones de la arquitectura deberán incorporar herramientas automatizadas de \textbf{Linaje de Datos (\textit{Data Lineage})} y observabilidad continua. Estas herramientas permitirán mapear visual e interactivamente las dependencias de los datos extremo a extremo (desde el origen transaccional hasta el panel analítico), habilitando alertas proactivas sobre caídas en la calidad del dato y reduciendo el tiempo de resolución de incidentes (MTTR) en la malla.


% =============================================================================
\section{Bibliografía}
% =============================================================================

\begin{thebibliography}{99}

\bibitem{dehghani2022}
Dehghani, Z. (2022). \textit{Data Mesh: Delivering Data-Driven Value at Scale}. O'Reilly Media.

\bibitem{machado2022}
Machado, I. N., Salgado, C., \& Maciel, P. (2022). Data Mesh: Concepts and Principles of a Paradigm Shift in Data Architectures. \textit{Procedia Computer Science}, 196, 263--271.

\bibitem{dixon2010}
Dixon, J. (2010). \textit{Pentaho, Hadoop, and Data Lakes}. James Dixon's Blog.

\bibitem{fowler2004}
Fowler, M. (2004). \textit{StranglerFigApplication}. MartinFowler.com. Recuperado de: \url{https://martinfowler.com/bliki/StranglerFigApplication.html}

\bibitem{brikman2019}
Brikman, Y. (2019). \textit{Terraform: Up \& Running: Writing Infrastructure as Code}. O'Reilly Media.

\bibitem{aws2023}
Amazon Web Services. (2023). \textit{AWS Well-Architected Framework: Data Analytics Lens}. AWS Whitepapers.

\bibitem{brewer2012}
Brewer, E. (2012). CAP twelve years later: How the ``rules'' have changed. \textit{Computer}, 45(2), 22-29.

\bibitem{abadi2012}
Abadi, D. (2012). Consistency tradeoffs in modern distributed database system design: CAP is only part of the story. \textit{Computer}, 45(2), 37-42.

\bibitem{finops2021}
FinOps Foundation. (2021). \textit{Cloud FinOps: Collaborative, Real-Time Cloud Financial Management}. O'Reilly Media.

\bibitem{nist2023}
Hu, V. C., Ferraiolo, D., \& Kuhn, R. (2006). \textit{Assessment of Access Control Systems}. National Institute of Standards and Technology (NIST).

\bibitem{aws_analytics_2021}
Amazon Web Services. (2021). \textit{Build a Data Mesh driven by Domain-Oriented Design}. AWS Architecture Center.

\bibitem{c4model}
Brown, S. (2018). \textit{Software Architecture for Developers: The C4 model for visualising software architecture}. Leanpub.

\end{thebibliography}

\end{document}